\chapter{Gemischte Paaroperationen}
\label{b14guide:ch:mixed}

\section{Allgemeingültige Ungleichungen}

\begin{theorem}[Die beiden Distributivitätsungleichungen]
\label{b14guide:thm:distributive-inequalities}
Existieren alle Paarwerte, die in der jeweiligen Zeile auftreten,
einschließlich der inneren Teilwerte, so gilt
\[
\begin{aligned}
 (x\wedge y)\vee(x\wedge z)&\leq x\wedge(y\vee z),\\
 x\vee(y\wedge z)&\leq(x\vee y)\wedge(x\vee z).
\end{aligned}
\]
Gleichheit folgt aus diesen Existenzbedingungen im allgemeinen
nicht. Es wird kein Distributivitätsaxiom vorausgesetzt.
\end{theorem}

\begin{theorem}[Gemischte Ungleichung bei vergleichbaren äußeren Operanden]
\label{b14guide:thm:modular-inequality}
Sei $x\leq z$. Sind $y\wedge z$, $x\vee(y\wedge z)$,
$x\vee y$ und $(x\vee y)\wedge z$ vorhanden, so gilt
\[
 x\vee(y\wedge z)\leq(x\vee y)\wedge z.
\]
Auch unter $x\leq z$ muß hier keine Gleichheit gelten.
Die Ungleichung verlangt kein Modularitätsaxiom.
\end{theorem}

\section{Gleichheiten bei vergleichbaren inneren Operanden}

\begin{theorem}[Distributive Gleichheiten für ein vergleichbares Paar]
\label{b14guide:thm:comparable-distribution}
Seien $y\leq z$ und $x\wedge y$, $x\wedge z$ vorhanden.
Dann existieren alle Paarwerte in
\[
 x\wedge(y\vee z)=(x\wedge y)\vee(x\wedge z).
\]
Seien $y\leq z$ und $x\vee y$, $x\vee z$ vorhanden.
Dann existieren alle Paarwerte in
\[
 x\vee(y\wedge z)=(x\vee y)\wedge(x\vee z).
\]
Dies sind lokale Folgerungen aus der Vergleichbarkeit
von $y,z$; eine Totalordnung auf $P$ wird nicht vorausgesetzt.
\end{theorem}